\begin{lemma}[Step 1, Subcase 2.1: the hierarchy on the recolored prefix]
Let $G,e_0,[k],T,\varphi,\mathcal E$ be as in the positive-rung setting,
with maximum property, and assume the complete six-part lower-rung package.
Let $H$ be a good hierarchy over a hierarchy base, and assume the large
estimate $|\overline\varphi(T_n)|\geq2n+11$ and elementarity of $V(T)$.

Let $\gamma\ne\beta$ and let $\mu$ be proper, leave exactly $e_0$
uncolored, and be $(T_n,D_n,\varphi)$-stable.  Suppose exact ordered
prefixes satisfy
\[
 H.\mathit{starred}(d)\preccurlyeq B^-\prec C\preccurlyeq T,
 \qquad C=B^-+g
\]
for some exact edge identity $g$.  Assume every exact edge incident with
$B^-$ has the same color under $\mu$ and $\varphi$, and every vertex of
$B^-$ has the same missing-color set under the two colorings.  At the last
vertex $b$ of $C$, assume $\gamma$ is missing under $\varphi$, $\beta$ is
missing under $\mu$, and an incident exact edge has color $\gamma$ under
$\mu$.  On every vertex of $C$, assume that membership in the missing set
is unchanged for colors other than $\gamma,\beta$.  Finally assume that
$\beta$ is not missing on $T$ under $\varphi$ and that $\gamma$ is absent
from the last indexed Gamma union of $H$.

Then $\mu$ is a \noderef{ModColoring} of $\mathcal E$, and $C$ carries a
\noderef{CorrespondingExtendedTashkinovTree} $\mathcal E_C$ under $\mu$
with maximum property, an actual \noderef{HierarchyBase}, and a genuine
good hierarchy $K_C$ with the same divider count as $H$.  The old terminal
PE route, or the exact old SE/RE seed followed by its route, transports to
a \noderef{TashkinovAugmentationSequence} ending at $C$.  The replacement
retains the large estimate, and both $\gamma$ and $\beta$ are absent from
the last indexed Gamma union of $K_C$.
\end{lemma}
\begin{proof}
Write $d=H.\mathit{dividerCount}$ and $A=H.\mathit{starred}(d)$.  The
lower field of \noderef{ExtendedTashkinovTree}, the PE and non-PE clauses
of \noderef{HierarchyBase}, the left-segment field of
\noderef{OrderedTreeUnion}, and the nesting clauses of
\noderef{TashkinovHierarchy} place the last ladder tree $T_n$ below $A$,
hence below $B^-$.  Therefore
\[
                    T_n\preccurlyeq A\preccurlyeq B^-
                       \prec C\preccurlyeq T.                 \tag{1}
\]

We first construct the exact terminal augmentation route to $C$.  Inspect
the last positive certificate in \noderef{TashkinovSeries} and
\noderef{TashkinovSeriesStep}.  Positivity excludes the initial
constructor.  In PE, the certificate contains a
\noderef{TashkinovClosure} of $T_n$ under $\varphi$.  In SE or RE it
contains an individual identity $f_n$, the exact equality $F_n=\{f_n\}$,
the exact seed $T_n+f_n$, and a closure of that seed.  The upper field of
\noderef{ExtendedTashkinovTree} puts $T$, and hence $C$, below the last
tree of this closure.

In PE, truncate the finite
\noderef{TashkinovAugmentationSequence} at the first member whose length
is $|C|$.  Such a member exists because every augmentation increases the
length by one, starts at $T_n$, and ends at a tree containing $C$ as an
exact prefix.  The member and $C$ are equal: both are exact initial
segments of the closure tree with equal lengths, so every vertex and
individual edge entry agrees.  In SE/RE, the seed and $C$ are exact
initial segments of the same closure tree.  Relation (1) gives
$|T_n|+1\leq |C|$, so entrywise comparison makes the exact seed, including
$f_n$, an initial segment of $C$.  The same truncation argument then gives
\[
 \begin{cases}
 T_n\longrightarrow C,&\Theta_n=\mathrm{PE},\\
 T_n+f_n\longrightarrow C,&\Theta_n\in\{\mathrm{SE},\mathrm{RE}\},
 \end{cases}                                                  \tag{2}
\]
under $\varphi$, retaining $F_n=\{f_n\}$ and the exact one-edge seed in
the second case.

Transport (2) step by step to $\mu$.  Before the final step into $C$,
the old prefix is contained in $B^-$.  Its new displayed edge is incident
with that old prefix, so the assumed incident-edge agreement fixes the
color of the exact identity.  The missing-color witness for the
augmentation is a vertex of the old prefix, so the assumed vertexwise
missing-set equality preserves it.  For the final step, its old prefix is
$B^-$ and its exact new edge is the identity $g$ in the stated
\noderef{TreeSequenceAdjoinEdge} certificate, so the same two arguments
apply.  Length, attachment, and exact initial-segment fields are
coloring-free.  These are all fields of
\noderef{TashkinovAugmentation}; induction over the finite sequence gives
the same alternatives (2) under $\mu$.  Denote their disjunction by
$\mathsf{Elig}_C$.

Apply \noderef{Td} to the given maximum property, positivity, complete
lower-rung package, properness, exact one-uncoloredness, and terminal
stability of $\mu$.  Its conclusion is precisely
\[
                         \noderef{ModColoring}(\mathcal E,\mu).\tag{3}
\]
Relation (1) supplies the exact strict extension of $T_n$, and
$\mathsf{Elig}_C$ supplies the appropriate PE or exact SE/RE alternative
in the universal clause of (3).  That clause yields an actual
\noderef{CorrespondingExtendedTashkinovTree} certificate
$\mathcal C_C$ on $C$.  By \noderef{MaximumPropertyTransport}, its
replacement ETT $\mathcal E_C$ has maximum property and the same positive
rung number $n$.

We next build the actual hierarchy base under $\mu$.  The correspondence
certificate identifies every ladder tree, raw stored set, exact
connecting-edge set, and extension kind.  If the terminal kind is not PE,
take the replacement core to be the common $T_n$, and copy the old
positive support position and the mathematically unused closure field.
The terminal missing-set compatibility is the vertex clause of the given
stability, and the nonparallel field of \noderef{HierarchyBase} is then
definitional.

Suppose the terminal kind is PE.  Extract the actual terminal PE
constructors of the original and replacement
\noderef{TashkinovSeriesStep}; write their maximum-defective data as
$(v,u,f,\delta_n,\gamma_n)$ and
$(v',u',f',\delta'_n,\gamma'_n)$.  Their singleton connecting-edge
equations, together with connecting-edge agreement in $\mathcal C_C$,
give
\[
                         \{f'\}=F'_n=F_n=\{f\},\qquad f'=f.    \tag{4}
\]
The endpoint fields of \noderef{MaximumDefectiveChoice} identify the
ends of this individual edge with both $\{v,u\}$ and $\{v',u'\}$.  The
two \noderef{DefectiveVertex} clauses put $v,v'$ in the common $T_n$,
while their other-end clauses put $u,u'$ outside it.  Thus the equality of
unordered endpoint pairs has only the orientation $v'=v$, $u'=u$.  The
old support position therefore records the actual replacement support.

Use the original core and closure $R_n$ for the replacement PE base.  We
do not assert that $R_n$ is an ordered prefix of the core.
The vertex-set field of \noderef{OrderedTreeUnion}, followed by the
hierarchy-core prefix and (1), gives only the inclusions that are needed:
\[
 V(R_n)\subseteq V(\mathit{core})\subseteq V(A)\subseteq V(B^-).\tag{5}
\]
Its left-segment field retains the common $T_n$ entries, and its
every-new-right-entry field retains the exact edge and vertex identities
of the added closure portion.

Transport the closure by incidence.  Every displayed edge of $R_n$,
including every new edge in its augmentation sequence, is incident with
its displayed new vertex in $V(R_n)\subseteq V(B^-)$.  Its exact color is
therefore fixed.  Every missing witness used by an augmentation also lies
in $V(R_n)$, so its missing set is fixed.  Hence every
\noderef{TashkinovAugmentation} in the closure remains valid under $\mu$.
If a boundary identity of $V(R_n)$ under $\mu$ had a color missing on
$R_n$, both its exact color and a missing witness would be unchanged on
$B^-$; it would be the same forbidden boundary identity under $\varphi$.
Thus the same $R_n$ is a genuine \noderef{TashkinovClosure} under $\mu$.
Together with (4), the retained support, and the coloring-free
\noderef{OrderedTreeUnion}, this proves every PE field of a replacement
base $B_C$.  This incidence argument, rather than a false prefix claim
about $R_n$, is the promised PE closure transport.

We construct the hierarchy on $C$.  If $A\prec B^-$, apply
\noderef{GoodHierarchyTerminalPrefixRepackage} to the original hierarchy
and the exact prefix $B^-$.  It yields the same $d$ nonterminal stages,
core, closure, stored sets, Gamma pairs, and a good hierarchy whose
terminal stage is $B^-$.  If $B^-=A$ and $d>0$, apply
\noderef{GoodHierarchyPrefixRepackage} at the last ordinary stage.  It
returns the good $(d-1)$-divider prefix and the separate last internal
minus-closed certificate; append $A$ as the $d$-th ordinary stage below
$C$.  If $B^-=A$ and $d=0$, take stage zero to be $T_n$, terminal stage
$C$, and core $A$ from $B_C$.  Relation (1) supplies every nesting and
strictness field.  In no case is a zero-rung ETT used.

Thus define $K_C$ over $B_C$ with exactly $d$ dividers, the old stages
through $A$, and terminal stage $C$.  Every starred stage lies in
$B^-$, where edge colors and vertex missing sets agree.  Correspondence
gives the same raw stored sets; expanding
\noderef{TashkinovStoredColors} and
\noderef{HierarchyStoredColorsAt} gives equality of every indexed stored
set.  Use the old Gamma family.  Pair cardinality, disjointness, and
unique refresh therefore copy.  All internal exact colored boundaries
are fixed, and the applicable prefix-repackaging result supplies the last
internal minus-closed clause.  In PE, the rebuilt base and the transported
closure above copy the parallel-base cardinality and non-minus-closed
clauses; in non-PE the implication is vacuous.

It remains to check terminal avoidance.  Expand
\noderef{HierarchySegmentColors} and the totalized
\noderef{HierarchyFirstMissingPosition}.  At indices below $d$, the next
stage lies in $B^-$ and all relevant data are unchanged.  At index $d$,
all displayed edges of $C$ retain their colors: old displayed edges are
incident with $B^-$, and the last identity $g$ is incident with $B^-$ by
the adjoin certificate.  For a color outside $\{\gamma,\beta\}$, the
assumed comparison of missing sets on $C$ and terminal prefixhood make
the new segment a subset of the old one.

For $\gamma$, elementarity of $V(T)$ and its old missingness at the last
vertex $b$ make $b$ its unique missing vertex in $T$.  No vertex of
$B^-$ therefore misses $\gamma$ under either coloring, and the stated
incident $\gamma$-edge under $\mu$, together with
\noderef{PresentColors} and \noderef{MissingColors}, says that $b$ does
not miss it under $\mu$.  Thus no vertex of $C$ misses $\gamma$, and the
old first position and new no-candidate fallback both cap the relevant
segment at the last position.  For $\beta$, the hypothesis that it is not
missing on $T$ says it is absent from every old missing set, while its
new missingness at $b$ makes $b$ the first new missing vertex.  The new
segment is again only the unchanged portion before the last position.
Old avoidance implies new avoidance in every case.  We have now checked
every field of a genuine \noderef{GoodHierarchy} witness $\mathcal G_C$.

The old exclusion of $\gamma$ transfers through the equality of indexed
stored sets and Gamma pairs.  If $\beta$ belonged to the last indexed
Gamma union, \noderef{GoodHierarchyReservedDisjoint} would make it
missing on $A$ and hence on $T$, contradicting the hypothesis.  Thus
$\beta$ is excluded as well.  Finally, correspondence identifies the
replacement last ladder tree with the exact old $T_n$, and the vertex
clause of stability identifies its missing-color set under $\mu$ with
the old one.  Taking cardinalities transports the large estimate.  The
objects $\mathcal C_C,B_C,K_C,\mathcal G_C$, (2), and (3) therefore
satisfy every conclusion.
\end{proof}
